\section{Introduction}\label{sec:intro}

\subsection{The problem}

Can a deterministic multitape Turing machine transpose a $d\times d$ matrix with $b(d)=\max(1,\clog d)$-bit entries,
written row by row, in time linear in the input length $L=d^2b(d)$? This is the statement~(T) of the first paper of this
series~\cite{FirstPaper2026}. By a reduction of Harvey and van der Hoeven~\cite[Cor.~6.9]{HvdH2025}, a linear-time
multiplication machine would give a linear-time transposer, so a proof of $\neg$(T) would give the first superlinear lower
bound for off-line integer multiplication on multitape Turing machines. The question is open, and this paper does not
decide it.

For oblivious machines the series has reduced $\neg$(T) to two statements. The fourth
paper~\cite[Theorem~4.10 and Corollary~4.11]{FourthPaper2026} shows that depth reduction of the time--space dependency
graph by a cheap genie, together with an index-coding lower bound on the side information of the remaining bounded-depth
cones (the \emph{crux}), rules out oblivious linear-time transposers. \emph{Directed Recovery Structures in Tape
Computations}~\cite{DirectedRecovery2026} shows that the cone relations of linear-time schedules have few coordinates
whose input and demand ports are close in the paired incidence graph, and that an affirmative answer to one machine-free
question would supply the crux on those cone relations~\cite[Question~9.1 and Proposition~9.2]{DirectedRecovery2026}: do
storage codes on bounded-degree recovery relations with this paired-distance separation have rate at most
$\frac12+o(1)$, for every alphabet?

This is a storage-code paper, like the series paper \emph{Covers and Local Generation in Linear Storage
Codes}~\cite{Covers2026}. We prove the exact half bound for all codes on large classes of recovery relations without short
paired paths, and show where the method stops. A recovery relation assigns to each coordinate $v$ the set $N(v)$ of
coordinates from which it must be recoverable. For a polymatroid $h$ valid for $N$, that is, with
$h(N(v)\cup\{v\})=h(N(v))$ for every $v$, the bound reads
\[
  2h(V)\;\le\;\sum_{v\in V}h(\{v\}).\tag{$\star$}
\]
Applied to the entropy of a uniformly random codeword, \eqref{eq:star} gives $\card C^2\le\card\alpha^{\card V}$ for
every storage code $C\subseteq\alpha^V$, linear or not, over every alphabet~\cite{FourthPaper2026}: no code beats the
pairing. Our hypotheses are exact rather than asymptotic. The central one is that the relation has no \emph{transitive
triangle} $x\to y\to z$, $x\to z$, where $u\to v$ means $u\in N(v)$; for loopless relations this says exactly that every
coordinate has paired distance at least five (\cref{lem:tt-paired}).

\subsection{Results}

\paragraph{The engine (\cref{sec:engine}).} \emph{Lemma~E} (nested elimination) eliminates a piece $H$ one vertex at a
time. Each vertex must be recovered from the vertices already eliminated, from one record other than its own, and possibly
from one vertex outside the piece; steps that use the same outside vertex must be nested. Then
$2h(H)\le\sum_ih(P_i)+\sum_{w\notin H}I(w;H)$ (\cref{lem:E}). Without outside vertices this is the inequality form of the
exact partition identity of the circular memo~\cite[Theorem~2.1]{CircularMemo2026}; the nested charging is new.
\emph{Theorem~D} (decomposition) contracts certified pieces one after another, each in the relation left by the earlier
ones, and proves \eqref{eq:star} whenever the vertex set has such a decomposition (\cref{thm:D}). Relative $2$-dicolourings
(Lemma~$\Bdir$, \cref{lem:Bdir}) and global certificates, which pull back along homomorphisms (\cref{lem:pullback}), are
the basic pieces.

\paragraph{The main theorem.} Every relation in the following classes has a decomposition, or a global certificate, and
hence satisfies \eqref{eq:star}; every storage code on it satisfies $\card C^2\le\card\alpha^{\card V}$.
\begin{enumerate}[label=(\alph*),leftmargin=2.5em]
\item $2$-dicolourable relations, in particular bipartite graphs (\cref{prop:twoAcyclic});
\item graphs of circular chromatic number below three, at any degree; this is the circular memo's
theorem~\cite[Corollary~3.2]{CircularMemo2026}, re-derived here as a direct instance (\cref{thm:circular});
\item triangle-free graphs of maximum degree three, and every relation whose arcs join adjacent vertices of such a graph
(\emph{Theorem~N3}, \cref{thm:N3});
\item loopless TT-free relations whose underlying graph has maximum degree three, directed triangles included
(\emph{Theorem~U3}, \cref{thm:U3});
\item loopless TT-free relations whose vertices with at least three parents induce an acyclic sub-relation, with no bound
on in- or out-degrees (\emph{Theorem~T}, \cref{thm:T}), and the version in which only parents inside a vertex's strongly
connected component count (\emph{Corollary~$\Tint$}, \cref{cor:Tint});
\item loopless relations partitioned into blocks with private outside parents, each block $K$ chordless subcubic or a
chorded digon path, with $N[K]$ TT-free, or internally certified (\emph{Theorem~PC}, \cref{thm:PC});
\item every TT-free Cayley relation $\Cay(G,\{s,s^{-1},t\})$, with arcs $g\to gx$, for every finite group $G$, when
$\{s,s^{-1},t\}$ has three elements and does not contain~$1$ (\emph{Corollary~CC}, \cref{cor:CC}).
\end{enumerate}
Corollary~$\Tint$ strictly contains the TT-free case of the two-input theorem of \emph{Directed Recovery
Structures}~\cite[Theorem~5.1]{DirectedRecovery2026}, and Theorem~U3 is the all-code version, on subcubic underlying
graphs, of the fourth paper's scalar linear Theorem~S~\cite[Theorem~6.13]{FourthPaper2026}. The certificates of
Theorem~PC come from \emph{Lemma~S} for chorded paths (\cref{lem:S}) and from sweeps of chorded circulant cycles
(\emph{Proposition~CC}, \cref{prop:CC}). The relation $\Cay(\Z_9,\{1,-1,3\})$ lies in none of the classes (a)--(e); it is covered
by~(g), and by~(f) with a single block.

\paragraph{Exact values (\cref{sec:consequences}).} On a triangle-free graph $G$ of maximum degree three, every storage code
satisfies $2\log\card C\le(\card V+\card S-\card{I(S)})\log\card\alpha$ for every vertex set $S$, where $I(S)$ is the set
of isolated vertices of $G-S$ (\cref{thm:barrier}). With the fractional Tutte--Berge formula (\cref{prop:fracTB}), the
guessing number equals the fractional matching number $\nustar(G)$, attained over every alphabet of square size
(\cref{thm:gn,cor:gn}). This proves, for these graphs and with the asymptotic guessing number taken either as the
supremum or as the limit over alphabet sizes, Christofides and Markström's conjecture that fractional clique covers are
optimal guessing strategies~\cite{ChristofidesMarkstrom2011,Dang2015}, which is false in
general~\cite{Baber2017,CameronDangRiis2016}; it was known before for perfect graphs, bipartite graphs in particular, and odd
cycles~\cite{ChristofidesMarkstrom2011}, and for prisms and certain chorded cycles~\cite{MMV2019}. It also proves Dang's
conjecture for planar triangle-free graphs at maximum degree three~\cite[Conjecture~15]{Dang2015}. Dually, the broadcast
rate of index coding is $\card V-\nustar(G)$ (\cref{thm:index,cor:broadcast}). Relations that become one of our classes
after deleting a set $Z$ have storage codes of rate at most
$(\card V+\card Z)/(2\card V)$ (\cref{cor:paid}).

\paragraph{Limits (\cref{sec:limits}).} The natural rules for choosing pieces fail: $\Cay(\Z_9,\{1,-1,3\})$ has no proper
piece with private outside parents (\cref{prop:no-private}); a $15$-vertex relation with in- and out-degrees at most three
has a unique minimum nonempty core, which is not a piece, although the relation has a global certificate
(\cref{rem:X15}); some chorded cycles have no valid sweep (\cref{rem:sweeps}); and the Petersen graph has no global
certificate, although Theorem~N3 gives it a decomposition (\cref{rem:petersen}). The half bound fails on the folded
cubes, whose odd girth equals their degree (\cref{rem:foldcube}, from a rank formula of Barg, Schwartz and
Yohananov~\cite{BSY2024}); so an exact half bound that assumes only $b_E(r)=0$, for a fixed odd $r\ge3$, cannot
hold at maximum degree $r+2$. For $r=3$ it fails already at maximum degree four: a $4$-regular triangle-free graph on
thirteen vertices, of odd girth five, carries a linear storage code over~$\F_4$ with $\card C=4^7$ (\cref{ex:deg4}), so
Theorem~N3 is sharp in the degree.

\paragraph{What this means for (T) (\cref{sec:construction}).} These are unconditional storage-code theorems; they do not
decide~(T). They constrain constructions: in the cone relations of a correct oblivious transposer left by
a genie of cost $\eps L$, no induced sub-relation on more than $2\eps L$ coordinates may lie in any of the classes above.
They also locate the remaining work for an impossibility proof: \cref{conj:P}, that every loopless TT-free relation of in-degree at most three has a decomposition,
would give the degree-three case of the paired-distance question (\cref{prop:dn3-transfer}).

\subsection{How to read this paper: the formal proofs}

Apart from three results of the fourth paper that we cite (\cref{thm:codes,lem:closure,lem:fewer-arcs}), which are
Lean-verified there, every numbered theorem, proposition, lemma, corollary and example of this paper is formalized in
Lean, including the guessing number and the broadcast rate (\cref{cor:gn,cor:broadcast}), the conditional transfer of
\cref{prop:dn3-transfer}, \cref{prop:no-private}, \cref{ex:D9} and \cref{ex:deg4}; the finite checks are evaluated by the
Lean kernel.
Not formalized are the remaining claims of the numerical remarks of \cref{sec:limits}: that the $15$-vertex relation
$\Ex$ is strongly connected and not $2$-dicolourable and that its minimum nonempty core is not a piece; that the twelve
chorded cycles of \cref{rem:sweeps} have no valid sweep; and that the Petersen graph has no global certificate. They are
reproducible with the scripts in \leanfile{tools/}. Not formalized either are the cited rank formula for folded cubes,
whose case of the Clebsch graph is Lean-verified in the fourth paper, the cited theorem of Atkins, Rombach and Skerman
after \cref{cor:gn}, and the heuristic \cref{rem:degree4}. The certificates printed in the remarks are Lean-verified.
None of the unformalized claims is used in a proof. Each numbered result states its status as described in \cref{sec:prelim-lean}.

\subsection{Related work}

\paragraph{Guessing numbers and graph entropy.} Guessing games and guessing numbers were introduced by
Riis~\cite{Riis2007}. Christofides and Markström~\cite{ChristofidesMarkstrom2011} studied undirected graphs, introduced
the fractional clique cover strategy and conjectured it to be optimal; the conjecture was disproved by Baber,
Christofides, Dang, Vaughan and Riis~\cite{Baber2017}, with triangle-free counterexamples by Cameron, Dang and
Riis~\cite{CameronDangRiis2016}. Dang's thesis~\cite{Dang2015} computes guessing numbers of several families of planar
triangle-free graphs; those of maximum degree three are special cases of Theorem~N3. Gadouleau and
Riis~\cite{GadouleauRiis2011} relate the guessing number of a digraph to the independence number of its guessing graph and
construct digraphs with large guessing numbers. The bounds there come from Shannon and non-Shannon inequalities or from
explicit strategies. Mazumdar, McGregor and Vorotnikova~\cite{MMV2019} bound storage capacity by covering a graph with
gadgets of known capacity, which already gives combinatorial certificates for Shannon-type upper bounds on some families.
Our certificates are of a different kind: elimination orders whose steps may use one outside vertex, with nested charges
(Lemma~E), composed by contraction rather than by covering (Theorem~D).

\paragraph{Storage codes and index coding.} Storage codes on graphs and their duality with index coding were studied by
Mazumdar~\cite{Mazumdar2014}, by Shanmugam and Dimakis~\cite{ShanmugamDimakis2014}, and by Arbabjolfaei and
Kim~\cite{ArbabjolfaeiKim2015}. Index coding was introduced by Birk and Kol~\cite{BirkKol2006} and formalized, with minrank
and the acyclic-subgraph lower bound, by Bar-Yossef, Birk, Jayram and Kol~\cite{BBJK2011}; Blasiak, Kleinberg and
Lubetzky~\cite{BKL2013} bound and approximate the broadcast rate by linear programming. Mazumdar, McGregor and Vorotnikova
bounded the storage capacity above by vertex covers and below by fractional matchings~\cite{MMV2019}; the lower bound is
Christofides and Markström's fractional clique cover strategy~\cite{ChristofidesMarkstrom2011} in storage language, and on
triangle-free subcubic graphs it is exact (\cref{thm:gn}). Without a degree bound, triangle-freeness does not force
rate~$\frac12$: Barg and Zémor~\cite{BargZemor2022} and Barg, Schwartz and Yohananov~\cite{BSY2024} construct storage codes
on triangle-free graphs of growing degree with rates approaching $\frac34$ and~$1$, and an earlier triangle-free family of
rate approaching~$1$ is due to Golovnev and Haviv~\cite{GolovnevHaviv2022}, as noted in~\cite{BSY2024}. The
bounded-degree examples among these have degree at least five, such as the Clebsch graph (\cref{rem:foldcube}); we are
not aware of an earlier triangle-free example of maximum degree four (\cref{ex:deg4}), and at degree three there is
none (Theorem~N3).

\paragraph{The series.} The fourth paper~\cite{FourthPaper2026} proved that rate bounds pull back along homomorphisms
(Theorem~H), the half bound for scalar linear codes on triangle-free subcubic side information with its matching form
(Theorem~S), and that the Clebsch graph has guessing number ten. The circular memo~\cite{CircularMemo2026} proved the
exact partition identity on which Lemma~E builds, and the half bound for circular chromatic number below three. The
subcubic memo~\cite{SubcubicMemo2026} proved half bounds for reversible and low-rank codes on subcubic supports and, for
all codes, on triangle-free subcubic graphs covered by four-cycles; Theorem~N3 contains the latter. \emph{Directed
Recovery Structures}~\cite{DirectedRecovery2026} proved the paired-distance bound for cone relations, the two-input theorem
and the hyperfinite half bound, and posed the all-code paired-distance question that motivates this paper. The covers
paper~\cite{Covers2026} studies cover structure and local generation of linear storage codes. The companion
paper~\cite{EighthPaper2026}, published together with this one, is about the depth side: its Theorem~F proves depth
reduction for head walks of bounded nesting load, so the crux alone, or an affirmative answer to Question~9.1 of
\emph{Directed Recovery Structures}, excludes oblivious linear-time transposers whose work heads enter every cell boundedly
often.

To our knowledge, Lemma~E, Theorems~N3, U3, T and~PC, and Corollary~CC are new. Special cases were known: the guessing
numbers of perfect graphs and odd cycles~\cite{ChristofidesMarkstrom2011}; the storage capacity of prisms, of certain
chorded cycles and of odd cycles~\cite{MMV2019}; the broadcast rates of cycles and of the cubic circulants, and numerically
that of the Petersen graph~\cite{BKL2013}; the families of planar triangle-free graphs in Dang's thesis~\cite{Dang2015};
the half bound on square-covered subcubic graphs~\cite{SubcubicMemo2026}; and the half bound for circular chromatic number
below three, which is the circular memo's theorem~\cite{CircularMemo2026}.

\subsection{Organization}

\Cref{sec:prelim} fixes notation. \Cref{sec:engine} proves Lemma~E and Theorem~D and the tools for building
certificates. \Cref{sec:undirected} treats $2$-dicolourable relations and symmetric side information, \cref{sec:directed}
the directed classes, and \cref{sec:consequences} the barrier bound, guessing numbers and index codes. \Cref{sec:limits}
shows where the method stops, \cref{sec:construction} draws the consequences for constructions, and \cref{sec:open} lists
open problems. The appendices index the Lean declarations and explain how to build and check the library.
